\subsection{An independent replication, and the one training variable stage A held fixed}
\label{sec:stageb}

Everything above comes from one campaign. Stage B is a second one, run separately with its own seeds
and its own run record, and it does two jobs (Table~\ref{tab:stageb}). Two of its cells repeat stage
A's protocol exactly at $d\in\{100,200\}$, which tests whether the effects replicate. The other two
double the training sparsity from $p=0.01$ to $p=0.02$, which is the one variable stage A held fixed
throughout and the obvious place for a result of this kind to be fragile.

\paragraph{The replication is stronger than the original.} At the matched post-ReLU comparison the
network leads the affine probe by \SbLfourHundredDiff{} of a sparsity level at $d=100$
([\SbLfourHundredCiLow, \SbLfourHundredCiHigh], \SbLfourHundredNetWins{} of \SbSeeds{} models) and
\SbLfourTwoHundredDiff{} at $d=200$ ([\SbLfourTwoHundredCiLow, \SbLfourTwoHundredCiHigh],
\SbLfourTwoHundredNetWins{} of \SbSeeds{}), against stage A's \PrimLfourHundredPostDiff{} and
\PrimLfourTwoHundredPostDiff{} in the same cells. The direction and the growth with width replicate;
the magnitude is roughly double, which is a reminder that a cell mean over ten or twenty seeds is not
a precise quantity even when its sign is stable. The untrained arm replicates too, and on the other
side: \SbRandHundredDiff{} and \SbRandTwoHundredDiff{} levels to the probe, unanimously.

Every $L^2$ cell in stage B has both decoders at $s_{95}=0$ --- \SbLtwoJointFailures{} models across
\SbLtwoCells{} cells --- which is the same joint failure stage A shows and is reported as such rather
than as agreement.

\paragraph{Doubling the training sparsity does not break the effect, and improves the frontier.}
At the higher sparsity the decoder difference shrinks but never turns negative:
\SbLoadLfourFiftyDiff{} at $d=50$, with all \SbSeeds{} seeds tied exactly as they are at that width
under $p=0.01$, and \SbLoadLfourTwoHundredDiff{} at $d=200$
([\SbLoadLfourTwoHundredCiLow, \SbLoadLfourTwoHundredCiHigh], \SbLoadLfourTwoHundredNetWins{} of
\SbSeeds{}). The frontier moves the other way and by more: at $d=200$ the exact $\kappa_{\min}$ of the
$L^4$ code rises from \SbKappaLfourTwoHundred{} at $p=0.01$ to \SbKappaLoadLfourTwoHundred{} at
$p=0.02$, while $L^2$ stays pinned at \SbKappaLoadLtwoTwoHundred{} --- a code trained on denser states
admits affine recovery at higher sparsity, and the loss still decides whether it admits any.

That last pair is the cleanest small result in the campaign, because the two arms differ in one
variable and move in opposite directions on it: the geometry the code presents improves with training
sparsity ($\kappa_{\min}$ up by a third), while the network's edge over a supervised affine fit of
that same code narrows. The two quantities are measuring different things --- what an affine rule can
guarantee in the worst case, against what a particular fitting procedure recovers on average --- and
it is useful to see them come apart, because a reader who takes them as interchangeable will draw the
wrong conclusion from either.

\begin{table}[t]
\centering
\caption{Stage B: a separate campaign with its own seeds and run record. The two $p=0.01$ rows per
arm repeat stage A's cells at $d\in\{100,200\}$; the two $p=0.02$ rows double the training sparsity.
``mean difference'' is network minus best post-ReLU affine probe in $s_{95}$, paired within seed, and
$\kappa_{\min}$ is the exact minimum over all $F$ features. Generated from
\texttt{results/e7\_stageB/derived/}.}
\label{tab:stageb}
\small
% Generated by scripts/make_tables.py from results/. Do not edit by hand.
\begin{tabular}{llcccccc}
\toprule
code & $d$ & $F/d$ & $p$ & mean difference & 95\% interval & net / tie / probe & $\kappa_{\min}$ \\
\midrule
$L^4$ & 100 & 4 & 0.01 & +0.70 & [+0.40, +1.00] & 7 / 3 / 0 & 4.3072 \\
$L^4$ & 200 & 4 & 0.01 & +1.00 & [+0.70, +1.30] & 9 / 1 / 0 & 5.6212 \\
$L^4$ & 50 & 2 & 0.02 & +0.00 & [+0.00, +0.00] & 0 / 10 / 0 & 4.5074 \\
$L^4$ & 200 & 2 & 0.02 & +0.40 & [+0.10, +0.70] & 4 / 6 / 0 & 7.6479 \\
$L^2$ & 100 & 4 & 0.01 & +0.00 & [+0.00, +0.00] & 0 / 10 / 0 & 1.0003 \\
$L^2$ & 200 & 4 & 0.01 & +0.00 & [+0.00, +0.00] & 0 / 10 / 0 & 1.0076 \\
$L^2$ & 50 & 2 & 0.02 & +0.00 & [+0.00, +0.00] & 0 / 10 / 0 & 1.0000 \\
$L^2$ & 200 & 2 & 0.02 & +0.00 & [+0.00, +0.00] & 0 / 10 / 0 & 1.0031 \\
untrained & 100 & 4 & 0.01 & -1.80 & [-2.00, -1.50] & 0 / 0 / 10 & 3.3999 \\
untrained & 200 & 4 & 0.01 & -2.00 & [-2.00, -2.00] & 0 / 0 / 10 & 4.7230 \\
\bottomrule
\end{tabular}

\end{table}
